%
%
%
%
%
\section{Computer-Assisted Parts, Data, Validation and Lean}\label{app:cert}
Every computer-assisted result of the paper consists of a written reduction, a theorem with a proof in the text, and an explicit finite predicate checked over an explicit finite set. Table~\ref{tab:ca} lists these items, followed by the exact computations that the paper quotes as parameters or as results of exhaustive searches in a stated scope. The data (check vectors, designs, window systems, tables, certificates) are archived with the paper, each item with a verifier that reads only the data files, never the search programs that produced them. Searches (random and exhaustive searches over check vectors, randomized searches for window systems, SAT solvers) only produced the data (Section~\ref{sec:conventions}).

\begin{table}[p]
\centering
\footnotesize
\caption{The computer-assisted items: the finite check, its scope, and the independent verification. Exact computations that are quoted but are not part of a proof are listed in Table~\ref{tab:checks}.}
\label{tab:ca}
\setlength{\tabcolsep}{3pt}
\begin{tabular}{@{}>{\raggedright\arraybackslash}p{0.2\textwidth}>{\raggedright\arraybackslash}p{0.31\textwidth}>{\raggedright\arraybackslash}p{0.24\textwidth}>{\raggedright\arraybackslash}p{0.19\textwidth}@{}}
\toprule
result & finite check & scope & verified by\\
\midrule
Tab.~\ref{tab:residue} (Thm.~\ref{thm:residue}, Cor.~\ref{cor:korobov}) & exact $\ljit(\Lambda_v)^2$ of the listed check vectors by enumeration of $\Lambda_v$ in a box; for the entries marked E, all check vectors $(x,-1)$ & the $18$ rows and the Korobov column & three programs (search, audit, and a box enumeration written from the definitions); one further brute force at $s=4$, $p=1009$\\
Thm.~\ref{thm:block}(c) below its proof range & exact $\LB(F)$ of $U^*_4$, $U^*_5$ (Lemma~\ref{lem:cert}, $j=1$, carry lengths capped at $s-1$) & every prime $7\le p\le109$ ($s=4$), $19\le p\le109$ ($s=5$) & two programs at every prime, written from Lemma~\ref{lem:cert}; a third at $7\le p\le31$, $p=101$ ($s=4$) and $37\le p\le101$ ($s=5$)\\
Prop.~\ref{prop:primepower} & exact $\beta_q(k)$ for every weight vector $k$ & the $13$ rows of Tab.~\ref{tab:primepower} & three programs\\
Thm.~\ref{thm:p7} & no vanishing combination in $\LB(\mathsf w)$ & $5$ classes of steps, $\le7^6$ combinations each & three programs (App.~\ref{app:p7}); Lean kernel\\
Thm.~\ref{thm:A3} & $\mathcal P(p_0)$, $p_0=19\,673\,600$ & one exact evaluation & second evaluation of $\mathcal P$; Lean kernel\\
Rem.~\ref{rem:A3range} (extends Thm.~\ref{thm:A3}) & hypotheses of Prop.~\ref{prop:steer} and Lem.~\ref{lem:W}, crosses computed in $\F_p$ & $1\,035\,877$ primes in $[3\cdot10^6,1.97\cdot10^7]$ & checker with its own exact window bound\\
thresholds after Thm.~\ref{thm:A}, $s=4,5$ (quoted, not part of its proof) & $\mathcal P(p_0)$ & one exact evaluation each & spot check above $p_0$; Lean kernel\\
Cor.~\ref{cor:HB17} & conditions of Thm.~\ref{thm:HB}, coverings as set identities & $715$ primes, $3677<p<10^4$ & second program; Lean kernel\\
Thms.~\ref{thm:B2}, \ref{thm:S3}, \ref{thm:K4}, \ref{thm:S4}, Rem.~\ref{rem:B3} & response sets (Prop.~\ref{prop:automaton}) and the minimum of Lem.~\ref{lem:cert} & Tab.~\ref{tab:frames} & second program; Lean kernel\\
Tab.~\ref{tab:practice}, certified bounds & generator matrices rebuilt from Def.~\ref{def:residue}, exact $\ljit(\Lambda_v)^2$ & $p=13$, $17$, $23$, $31$ & three programs\\
\bottomrule
\end{tabular}
\end{table}

\begin{table}[tp]
\centering
\footnotesize
\caption{Exact computations that are quoted in the paper but are not part of a proof of a theorem, proposition or corollary (continuation of Table~\ref{tab:ca}).}
\label{tab:checks}
\setlength{\tabcolsep}{3pt}
\begin{tabular}{@{}>{\raggedright\arraybackslash}p{0.2\textwidth}>{\raggedright\arraybackslash}p{0.31\textwidth}>{\raggedright\arraybackslash}p{0.24\textwidth}>{\raggedright\arraybackslash}p{0.19\textwidth}@{}}
\toprule
result & finite check & scope & verified by\\
\midrule
Tab.~\ref{tab:sobolpar} (parameters of the close pairs) & orders $\nu$, $k$ and the levels $L$, $m$ & $s=2,\dots,9$ & second program\\
Cor.~\ref{cor:halton}, explicit pairs & the two points of each pair are closer than the bound & $395$ cases of bases and carrying coordinate & exact distances of the pairs\\
Rem.~\ref{rem:skelsearch} (exhaustive in its scope) & branch and bound over all check-vector sequences of the best class & $s=2$; $p=11$, $13$, $17$; $K_{\max}=10$, $8$, $11$ & independent reproduction\\
example after Cor.~\ref{cor:nodesign} & dimension of a space of relations & $s=4$, $J=3$; $p=11$, $K=4,\dots,11$; $p=13,17$, $K=4$ & two methods (and the audit)\\
Lem.~\ref{lem:witness} brackets & explicit close pairs of points & $K=18$ ($\F_4$), $K=75$, $64$ & second program for Thm.~\ref{thm:K4} only\\
\bottomrule
\end{tabular}
\end{table}

\subsection{Verification}\label{app:verification}
Every finite check used in the proof of a theorem, proposition or corollary was carried out by at least two programs, as listed in Table~\ref{tab:ca}, the second written from the statements of the paper only. Remark~\ref{rem:A3range} and the brackets of Appendix~\ref{app:sharp} for Theorems~\ref{thm:S3} and~\ref{thm:S4} rest on one program each; the thresholds for $s=4,5$ quoted after Theorem~\ref{thm:A} rest on one program and the Lean kernel.

\emph{Residue and block-periodic sequences.} The exact values $\ljit(\Lambda_v)^2$ of Table~\ref{tab:residue} were computed by the enumeration program of the search, recomputed by an independent enumeration in the audit, and checked a third time by a short program that enumerates, for a claimed value $L$, the box $|b_c|\le1+\lfloor\sqrt L\rfloor$ (every lattice vector with $\|(|b|-\mathbf 1)_+\|_2^2\le L$ lies in it); the same program checks the Korobov column and the four check vectors of Table~\ref{tab:practice}. The entry $s=4$, $p=1009$ was in addition confirmed by a brute force over the full box $|b_c|\le6$. For Theorem~\ref{thm:block}(c) both programs evaluate the one-digit relaxation from its definition, one from the two possible differences of positions, the other by the minimum over all positions $R$; both read the designs $U^*_4$, $U^*_5$ from the same data file, whose $s=4$ entries are the matrices printed after Theorem~\ref{thm:block}, and the second rejects a design whose largest diagonal entry is changed by one. The values $\beta_q$ of Proposition~\ref{prop:primepower} were computed by the original program, by the audit's field arithmetic over every power-basis encoding, and by a third program for the encodings of Table~\ref{tab:primepower}; the generator matrices of the four residue sequences of Table~\ref{tab:practice} were rebuilt entry by entry from Definition~\ref{def:residue}, and their values of $\mathcal J$ were recomputed by an independent pipeline.

\emph{Window systems.} The verifiers of the window systems were in addition run by an adversarial audit with its own implementation. The common verifier reads a record and checks: the constants $\lambda_c$ and the places are admissible modulo $p$; the crosses $\kappa_{c\leftarrow c'}$ computed from the $\lambda_c$ by Proposition~\ref{prop:LI}(c), in two equivalent normalisations, equal the stored ones; for a rational system, the integrality of the crosses; condition (i) of Proposition~\ref{prop:steer}; condition (ii) for every column, exactly (by an exact scan for at most two restricted crosses, by enumeration of $\F_p^j$ when $p^{j-1}$ is at most a threshold, $2\cdot10^6$ by default) or else by Lemma~\ref{lem:ws} in the stored order of the rows; and the exact window bound $\LB(W)$ from its definition. Deliberately corrupted records (a cross changed, an own window shortened) are rejected. In the last full run all $6728$ stored records passed, including records of the previous version that this paper no longer uses. The exact separations quoted for figures and tables come from exact minimum-distance programs on the generated points.

\subsection{Archive}\label{app:archive}
The archive accompanying the paper, the directory \texttt{archive} of the repository named on the title page (release v1.0.0, archived at Zenodo), has one directory per item:
%
%
\texttt{verify} (the common verifier of the window systems);
\texttt{residue\_tables} (the check vectors of Table~\ref{tab:residue} with the exact $\ljit(\Lambda_v)^2$ and the three enumeration programs);
\texttt{block\_lb} (the designs $U^*_s$ of Theorem~\ref{thm:block}(c) and the per-prime evaluations of $\LB(F)$ for $s=4,5$, with both programs);
\texttt{primepower} (the polynomials, weights and values $\beta_q$ of Table~\ref{tab:primepower}, with the three programs);
\texttt{p1\_certificates}, \texttt{mrA\_designs} and \texttt{lf\_designs} (the ordered window systems of an earlier version of this work, not used here);
\texttt{p7\_\allowbreak graph\_\allowbreak certificate} (Theorem~\ref{thm:p7});
\texttt{bs\_ruleB\_s3} (Theorem~\ref{thm:A3}, Remark~\ref{rem:A3range});
\texttt{bs\_\allowbreak s45\_\allowbreak thresholds} (Theorem~\ref{thm:A});
\texttt{selection\_rule} (the rule of Remark~\ref{rem:steerrule} and the window systems of the sequences in the figures);
\texttt{practice} (the four residue sequences of Table~\ref{tab:practice}, given by places, scale and check vector, with their generator matrices, and the ordered window systems of the comparison in Section~\ref{sec:practice}, with a stand-alone point generator);
\texttt{b2\_dropin} (the generator matrices of Theorem~\ref{thm:B2} as a drop-in table of $32$-bit direction numbers, with a reference implementation in C);
\texttt{failure} (the parameters of Table~\ref{tab:sobolpar} and the check of the explicit pairs of Section~\ref{sec:fail});
\texttt{exact\_relations} (the exhaustive search of Remark~\ref{rem:skelsearch} and the two computations of the example after Corollary~\ref{cor:nodesign}).
Each directory has a README with the statement it supports and the command that re-runs its check; for instance \texttt{python3 verify/verify\_designs.py} with a record file of the directory. Every check runs in less than half a minute on one core (the $s=5$ window systems in chunks). The requirements are Python~3, a C\texttt{++}17 compiler and, for the exact rational evaluations, SymPy. The record format of the window systems is described in \texttt{verify/README}. The automata of Section~\ref{sec:frames} and the range of Corollary~\ref{cor:HB17} are part of the Lean development below.

\subsection{Lean}\label{app:lean}
The written proofs are the primary evidence; the formalization is an additional, independent check of them, not a replacement. The library \texttt{GeoSeq} (Lean~4 with Mathlib), which extends the Lean development of \cite{Gru26geo}, contains formal statements and proofs of most results of this paper. Every formalized statement is restated in a separate statement file and compared with its proof by the Comparator tool, which re-checks the proofs with two independent kernels (the Lean kernel and NanoDa) and admits only the three standard axioms of Lean's classical logic (propositional extensionality, quotient soundness and choice). No proof uses \texttt{sorry} or \texttt{native\_decide}; the finite computations inside the formal proofs are evaluated by the Lean kernel itself.

\emph{Formalized results of this version.} Section~\ref{sec:fail} for prime bases: Definition~\ref{def:radinv}, Lemma~\ref{lem:carrypair}, Theorem~\ref{thm:nonqu}, Corollaries~\ref{cor:halton}, \ref{cor:faure}, \ref{cor:sobol}, \ref{cor:nied} and~\ref{cor:unscrambled}, Lemmas~\ref{lem:sobol} and~\ref{lem:nied}, Remark~\ref{rem:faurerate}, and the reorderings of Remark~\ref{rem:randomize}. Lemma~\ref{lem:resform}(a) and~(c), and Lemma~\ref{lem:prefix} for general level families together with Remark~\ref{rem:prefixcases}. Section~\ref{sec:residue}: Definition~\ref{def:residue}, Lemmas~\ref{lem:residue-t0} and~\ref{lem:residue-rel}, Theorem~\ref{thm:residue}, Corollaries~\ref{cor:korobov} and~\ref{cor:residueQU}, Definition~\ref{def:residuestar}, Lemmas~\ref{lem:residuestar-t0} and~\ref{lem:residuestar-rel}, Theorem~\ref{thm:residuestar}, Corollary~\ref{cor:dyadic}, Example~\ref{ex:brickwall}, and Theorem~\ref{thm:WF} for the definition of the brick wall used here (Definition~\ref{def:WF}). Section~\ref{sec:block}: Theorem~\ref{thm:block}(i)--(iii), its cases (a) and (b), and case (c) for $s=4$ and $s=5$ with the computer-assisted ranges, Theorem~\ref{thm:blocktwo}, Proposition~\ref{prop:implied} and Corollary~\ref{cor:gauss}. Section~\ref{sec:count}: Definition~\ref{def:exactrel}, Lemma~\ref{lem:reltype} (including the dimension of the spaces of relations), Theorem~\ref{thm:count}, Corollaries~\ref{cor:layer}, \ref{cor:N1}, \ref{cor:residuerows}, \ref{cor:nodesign} and~\ref{cor:density}, Propositions~\ref{prop:onerel} and~\ref{prop:blockrel}, and Lemma~\ref{lem:rigid}. Section~\ref{sec:joint}: Lemma~\ref{lem:corrS}, Theorem~\ref{thm:corrected} and Proposition~\ref{prop:joint}; Lemma~\ref{lem:shift}(b). Theorem~\ref{thm:coverage}, without its prime-power rows. Of the results kept from the previous version, the library covers, with statements that agree up to wording: Lemmas~\ref{lem:t0} and~\ref{lem:cert}, Remark~\ref{rem:onedigit}, Proposition~\ref{prop:law}, Example~\ref{ex:running}; Propositions~\ref{prop:LI}, \ref{prop:depthdir} and~\ref{prop:steer}, Lemmas~\ref{lem:volume} and~\ref{lem:antipode}, Theorem~\ref{thm:BS}, Lemmas~\ref{lem:W}, \ref{lem:D} and~\ref{lem:R}, Proposition~\ref{prop:P0}, Theorems~\ref{thm:A3}, \ref{thm:A} and~\ref{thm:HB}, Corollary~\ref{cor:HB17}, Lemma~\ref{lem:ws}; Theorems~\ref{thm:B2}, \ref{thm:S3}, \ref{thm:K4}, \ref{thm:S4} and~\ref{thm:p7} and Remark~\ref{rem:B3} with their automata (for the instances used), Lemma~\ref{lem:octave}, Corollary~\ref{cor:allprimes}, Lemma~\ref{lem:shift}(a), and Propositions~\ref{prop:N2} and~\ref{prop:rank2}.

\emph{Where the formal statements differ in form.} The results of Section~\ref{sec:fail} are formalized for prime bases $q=p$ only. Self-similar scrambles enter the formal proofs through their residue form, that is, they are defined by local multipliers as on the right-hand side of Lemma~\ref{lem:resform}(b); the equivalence with Definition~\ref{def:selfsim} is not formalized. The dimension bound of Proposition~\ref{prop:onerel}(i) and Proposition~\ref{prop:blockrel} are derived from Theorem~\ref{thm:count}, which makes them hold for every scramble of a frame, as the remark before Proposition~\ref{prop:onerel} indicates. Corollaries~\ref{cor:layer}(a) and~\ref{cor:density} are formalized in their form for one admissible range of levels; the statements about upper densities, which follow by covering the levels with admissible ranges, are not. For Theorem~\ref{thm:block}(c) the computer-assisted part is a checker of the one-digit relaxation, proved sound once and evaluated by the kernel at every prime $7\le p\le109$ ($s=4$) and $19\le p\le109$ ($s=5$); it shows $\LB(F)\ge1-(a^*_s+1)/p$, which is all the theorem needs, while the equality of $\LB(F)$ with this value is not formalized. The design multipliers of Definition~\ref{def:block} enter as explicit polynomials for the cases (a)--(c) and Corollary~\ref{cor:gauss}, with the system (a) and the orders of vanishing verified for every admissible prime; their existence for general designs is not formalized. Proposition~\ref{prop:joint}(ii) is formalized for one factor $\varphi_c(1+\varepsilon w^j)$; the reduction of a general multiplier to products of such factors is not. The square-lattice clause of Corollary~\ref{cor:gauss} is formalized as the inclusion of the occupied cells in the lattice. Theorem~\ref{thm:coverage} is formalized as an assembly of the cited results.

\emph{Not formalized.} Lemma~\ref{lem:resform}(b) and Lemma~\ref{lem:reschange}: the change of variable in characteristic $p$ needs the universal-coefficient argument of its proof, and the theorems that use Lemma~\ref{lem:resform}(b) are formalized for scrambles given in residue form instead. Proposition~\ref{prop:primepower} beyond $t=0$ and the relation (which hold over every finite field in the library), and hence the prime-power rows of Theorem~\ref{thm:coverage}: the values $\beta_q(k)$ would need explicit models of $\F_q$ with the power-basis encoding. Theorem~\ref{thm:block}(c) for $s\ge6$, which needs the reachability of the general design $U^*_s$ at the moved places. The prime powers in Section~\ref{sec:fail}, and the digital shifts and index scrambles of Remark~\ref{rem:randomize}. The computer-assisted tables (Tables~\ref{tab:residue}, \ref{tab:sobolpar}, \ref{tab:primepower} and~\ref{tab:practice}) and Remark~\ref{rem:A3range}, Lemmas~\ref{lem:riordan} and~\ref{lem:witness} (neither is used in the formal proofs), the binary-sequence clauses of Theorems~\ref{thm:K4} and~\ref{thm:S4}, Remark~\ref{rem:toeplitz}, and the remaining remarks and the figures.

\emph{Configurations.} The results new in this version are compared in $32$ Comparator configurations with $311$ theorems; the $22$ kernel evaluations for Theorem~\ref{thm:block}(c) with $s=5$ form one configuration per prime, for run time. All of them, and the Palomar entry below, were re-run once in sequence on the final state of the library and accepted by both kernels. The results kept from the previous version are compared in the configurations of that version.

\emph{Palomar entry.} A self-contained entry in the format of the Palomar collection of formalizations (one file of $195$ lines that imports only Mathlib) states the two headline results with their own definitions: no Halton-type sequence in pairwise coprime polynomial bases is quasi-uniform (Corollary~\ref{cor:halton}), and the residue sequence with the Korobov check vector and Faure places is a digital $(0,s)$-sequence for every prime $p>2^s$ with the separation bound of Corollary~\ref{cor:korobov} at every prefix $P_{sK-1}$ and bounded mesh ratio (Theorem~\ref{thm:residue}, Corollaries~\ref{cor:korobov} and~\ref{cor:residueQU}). Its definitions are checked against those of the library by bridge theorems.

\subsection{The Small-Base Bounds: Explicit Close Pairs}\label{app:sharp}
The certified bounds of Theorems~\ref{thm:S3}, \ref{thm:K4} and~\ref{thm:S4} are far below the exact separations of Table~\ref{tab:frames}, but for Theorem~\ref{thm:K4} (two programs), and for Theorems~\ref{thm:S3} and~\ref{thm:S4} by a single implementation, this is not a weakness of the bound. A carry of length $l$ exists only from level $K=l$ on, and the binding steps involve long carries, so the computable prefixes are too short to contain them. The next lemma turns a binding configuration into an actual pair of points.

\begin{lemma}[Explicit close pairs]\label{lem:witness}
Let the offset map be the same at every level, and fix $j$, a step $e$, carry lengths and carry digits of the moving coordinates, with the changes $\sigma_c$ of Lemma~\ref{lem:cert}. At a level $K$ at least the largest carry length, the cells with this carry pattern are given by prescribing some of their finest digits. If the affine map from the remaining cell digits to the geometric digits $(R_c)_c\in\F_q^{sj}$ is onto, then $\sep{P_{sK}}\le\bigl(\sum_c(f_j(e_c,\sigma_c)+2q^{-j})^2\bigr)^{1/2}$.
\end{lemma}

\begin{proof}
For each $c$ choose $R^*_c$ attaining the minimum in $f_j(e_c,\sigma_c)$, and, by surjectivity, a cell $n$ with the carry pattern and $R(n)=R^*$. The neighbour $n+e$ has the geometric digits $R^*_c\oplus\sigma_c$, so by (\ref{eq:cellunits}) the two points differ in coordinate $c$ by at most $f_j(e_c,\sigma_c)+2q^{-j}$ cell units.
\end{proof}

For Theorem~\ref{thm:K4} the binding configuration (carry lengths $9$ and $16$) is realisable at $K=18$, where the map has full rank; solving for the free digits gives two points of $P_{54}$ at distance $0.112940$ cell units. Hence $0.112930\le\inf_K\sep{P_{3K}}\le0.112940$ (computer-assisted): the constant of Theorem~\ref{thm:K4} is determined to $0.01\%$. The same construction gives $\inf_K\sep{P_{3K}}\le\sep{P_{225}}\le0.05616$ for Theorem~\ref{thm:S3} and $\inf_K\sep{P_{4K}}\le\sep{P_{256}}\le0.03176$ for Theorem~\ref{thm:S4}. So larger constants in these frames need other multipliers, not better bounds.

As binary sequences, the prefixes of $2^M$ points of Theorem~\ref{thm:K4} have base-$2$ quality parameter at most $3$ for $M\le18$, with $t=0$ at every full level $M=6K$; those of Theorem~\ref{thm:S4} have at most $4$ for $M\le20$ (computations), against $t=1$ and $t=3$ for the three- and four-dimensional Sobol' sequences.

\begin{remark}[Why these bases]\label{rem:height}
Heuristically: in a Kummer frame $\coef{n}\tau^{-b}=\prod_i\binom b{n_i}$ over the base-$q$ digits $n_i$ of $n$. For $q=3$ only $b\le1$ occurs, and for $q=4$, $(1+x)^2=1+x^2$ keeps all coefficients in $\{0,1\}$, so the automata stay small; over $\F_5$ the coefficients $2$ and $3$ appear, and the symmetric scrambles we screened had more than $300$ response columns per coordinate at $j=3$ and no positive bound.
\end{remark}
